\section{Objetivos}

\subsection{Objetivo general}

Implementar y administrar remotamente un servicio WordPress en arquitectura de dos capas sobre un hipervisor Tipo 1, aplicando controles de red, identidad, permisos de archivos y cifrado TLS verificables.

\subsection{Objetivos específicos alcanzados}

\begin{itemize}
  \item Instalar Proxmox VE en hardware físico y ejecutar tres VM simultáneamente.
  \item Instalar Ubuntu Server 26.04.1 LTS sin entorno gráfico en las VM web y DB.
  \item Utilizar openSUSE Leap 16.0 con KDE como estación de administración.
  \item Separar Apache, PHP y WordPress de MySQL mediante la red interna \texttt{10.50.5.0/24}.
  \item Cambiar SSH a \texttt{22222/TCP} y MySQL a \texttt{13306/TCP}.
  \item Restringir MySQL a la VM web y demostrar el bloqueo desde el cliente.
  \item Aplicar propiedad \texttt{wpadmin:www-data} y modos 750/640 en WordPress.
  \item Validar SFTP con FileZilla y una prueba negativa de permisos.
  \item Configurar UFW con política entrante denegada y reglas limitadas por origen.
  \item Emitir y renovar un certificado válido de Let's Encrypt para \texttt{laboratorio5.com}.
  \item Documentar resultados, incidencias y evidencia reproducible.
\end{itemize}
